\chapter{GİRİŞ}
\label{chap:giris}

Elektrik enerjisi sistemleri; artan talep, fosil kaynakların çevresel etkileri ve enerji arzının karbonsuzlaştırılması gereksinimi nedeniyle önemli bir dönüşüm geçirmektedir. Güneş ve rüzgâr gibi yenilenebilir kaynakların kurulu güç içindeki payı büyürken elektrik üretimi, merkezi ve öngörülebilir bir yapıdan dağıtık ve zamana bağlı bir yapıya doğru ilerlemektedir. Söz konusu dönüşüm, temiz enerji üretimini artırmakla birlikte üretim ile tüketimin her an dengelenmesi zorunluluğunu ortadan kaldırmamaktadır. Işınım ve rüzgâr koşullarına bağlı üretim değişimleri; güç dengesinin korunması, gerilim ve frekansın desteklenmesi ve yenilenebilir enerjinin etkin biçimde kullanılabilmesi için esnek kaynaklara duyulan gereksinimi artırmaktadır \citep{Chen2009,Aneke2016,IEAPVPS2024}.

Enerji depolama sistemleri (EDS), enerjiyi üretildiği anda tüketme zorunluluğunu azaltarak yenilenebilir kaynakların şebekeye entegrasyonunda bu esnekliği sağlayan başlıca yapılardan biridir. Bir EDS yalnızca batarya grubundan ibaret değildir; batarya, güç dönüşüm sistemi, batarya ve enerji yönetim birimleri, koruma, haberleşme ve ısıl yönetim alt sistemlerinin birlikte çalıştığı bütünleşik bir sistemdir \citep{Hannan2017,Plett2015}. Dolayısıyla sistem başarısı yalnızca hücrenin enerji yoğunluğuyla değil, batarya ile DA bara, yük ve şebeke arasındaki enerji akışını yöneten güç elektroniği arayüzünün verimi, güvenilirliği, hacmi ve maliyetiyle de belirlenir. Güncel maliyet çalışmalarında batarya paketi en büyük kalem olarak öne çıksa da çift yönlü güç dönüştürücüleri ve bunların yardımcı donanımları da sistem düzeyindeki maliyet ve kayıpların önemli bir parçasıdır \citep{Ramasamy2023,Lazard2024}. Bu tablo, hem bataryanın güvenli çalışma bölgesini koruyan hem de aynı yarı iletken donanımı farklı enerji akış yönlerinde daha etkin kullanan dönüştürücü mimarilerini önemli kılmaktadır.

Depolama arayüzünün performansı, bağlı olduğu yenilenebilir üretim sisteminin gerçek çalışma profiliyle birlikte değerlendirilmelidir. Fotovoltaik (FV) üretim, ışınım ve sıcaklık nedeniyle gün boyunca sürekli değişir; evirici yalnızca sınırlı sürelerle anma gücüne yakın çalışır ve çalışma ömrünün kayda değer bir bölümünü kısmi yük bölgelerinde geçirir. Dolayısıyla yalnızca anma noktasında verilen tepe verim değeri, yıllık enerji dönüşüm başarısını tek başına açıklamaz. Avrupa ve CEC ağırlıklı verim tanımları ile bölgesel ışınım dağılımlarına dayalı ölçütler, farklı güç seviyelerindeki verimlerin gerçek işletme açısından birlikte ele alınması gerektiğini göstermektedir \citep{Meinhardt2000,Sera2008,Sandia2015,Ongun2013}. Kısmi yükte sabit sürme, anahtarlama ve manyetik eleman kayıplarının aktarılan güce oranı büyürken akım dalga biçimi de anahtar doğrusal olmayanlıklarına, ölü zamana ve şebeke gerilimindeki bozulmalara daha duyarlı hâle gelir. Sonuçta düşük güç bölgesi, verim ile akım toplam harmonik bozulması (Total Harmonic Distortion, THD) arasındaki ödünleşimin belirginleştiği bir çalışma alanına dönüşür \citep{Kjaer2005,Chicco2009}.

Şebekeye bağlı dönüştürücülerde bu ödünleşim; topoloji, modülasyon, filtre ve kontrol yapısının ortak sonucudur. İki seviyeli eviricilere alternatif olan çok seviyeli topolojiler, çıkış gerilimini daha fazla seviyede oluşturarak yarı iletken gerilim zorlanmasını ve anahtarlama dalgalanmasını azaltabilmektedir. Üç seviyeli T-tipi evirici, nötr noktaya bağlanan çift yönlü anahtar kolu sayesinde $+V_{\mathrm{dc}}/2$, $0$ ve $-V_{\mathrm{dc}}/2$ gerilim seviyelerini üretir. T-tipi yapı; uygun eleman seçimi ve modülasyonla yüksek güç yoğunluğu ve yüksek verim potansiyeli sunmakta, aynı zamanda ticari yarı iletken modüllerle gerçekleştirilebilmektedir \citep{Wang2020}. T-tipi topolojinin şebekeye bağlı gerçek bir donanım üzerindeki uygulanabilirliği ve elektromanyetik uyumluluk gereksinimleri de deneysel çalışmalarla gösterilmiştir \citep{Hizarci2021}.

Bununla birlikte, topolojinin anma gücündeki üstünlüğü düşük güçte aynı ölçüde korunmayabilir. Çıkış akımının küçük olduğu bölgelerde T-tipi faz kolunun sıfır gerilim durumunu oluşturan nötr anahtarları, akımın işaret değiştirmesine izin veren senkron bir alçaltıcı dönüştürücü (synchronous buck) davranışı gösterebilir. Endüktör akımı bir anahtarlama periyodu içinde sıfıra ulaştıktan sonra nötr anahtar iletimde tutulduğunda akım ters yönde büyüyebilir; böylece şebekeye yararlı güç taşımayan dolaşım akımı, yarı iletken ve pasif elemanların RMS akımını ve iletim kayıplarını artırabilir. Süreksiz iletim ve diyot benzetimi (diode emulation) yaklaşımlarının hafif yükte ters akımı engelleyerek kayıpları azaltabilmesi, senkron alçaltıcı dönüştürücüler için bilinen bir ilkedir \citep{Erickson2001,Mohan2003,Chen2010}. Ancak bu ilkenin şebekeye bağlı üç seviyeli T-tipi eviricinin sıfır durumuna uygulanması, yalnızca bir anahtarın kapatılması değildir; akım yolunu, çıkış gerilimini, modülasyon bütünlüğünü ve harmonik içeriği birlikte değiştiren topolojiye özgü bir problemdir.

Tez kapsamında söz konusu problem için, düşük güç bölgesinde ters akım dolaşımını sınırlandırmak amacıyla nötr anahtarın belirli koşullarda susturulmasına dayanan yeni bir anahtarlama yöntemi ele alınmaktadır. Yöntem, dolaşım akımını ve buna bağlı kayıpları azaltmayı hedeflerken T-tipi eviricinin ürettiği gerilim darbelerini de değiştirmektedir. Verimdeki olası iyileşmenin karşılığında, özellikle üçüncü harmonik bileşeninde artış ortaya çıkabilmektedir. Ölü zaman, anahtar geçişleri ve parazitik kapasiteler gibi gerçek devre etkilerinin akım harmoniklerini çalışma noktasına bağlı biçimde değiştirebildiği bilinmektedir \citep{Zhang2014Deadtime}. Önerilen anahtarlama yönteminin başarısı bu yüzden yalnızca verim kazancıyla değil, aynı çalışma noktalarında THD ve şebeke akımı kalitesi korunarak değerlendirilmelidir.

Anahtarlama değişikliğinin oluşturduğu üçüncü harmonik bileşenini azaltmak için tez kapsamında Sin\textsuperscript{3} tabanlı bir ileri besleme kompanzasyonu (feedforward compensation) kullanılmaktadır. $\sin^3(\theta)=[3\sin(\theta)-\sin(3\theta)]/4$ özdeşliği, ölçülen şebeke geriliminin faz bilgisinden temel ve üçüncü harmonikle ilişkili anlık bir şekillendirme terimi üretilebilmesini sağlar. Üretilen terim, akım hatasından türetilen bağımsız bir geri besleme kontrolcüsü değil, mevcut akım kontrol döngüsünün modülasyon girişine eklenen bir kompanzasyon bileşenidir. Söz konusu terminolojik ayrım önemlidir: şebeke gerilimi ileri besleme yolları akım bozulmasını azaltabilse de dijital gecikme, LCL filtresi ve şebeke empedansı nedeniyle sistem dinamiğiyle etkileşebilir \citep{Lu2016,Wang2024}. Sin\textsuperscript{3} teriminin ideal doğrusal modelde karakteristik paydayı değiştirmemesi bu yüzden gerçek dönüştürücüde kararlılık incelemesini gereksiz kılmaz; kapalı çevrim modelinin, duyarlılık fonksiyonlarının ve düşük ile anma gücü çalışma noktalarının birlikte incelenmesi gerekir.

Tezin ikinci daralma ekseni, aynı T-tipi donanımın yalnızca evirici olarak değil, EDS'nin farklı enerji akış kiplerinde yeniden kullanılabilmesidir. Geleneksel bir depolama sisteminde şebeke eviricisi, şarj doğrultucusu veya güç faktörü düzeltme (Power Factor Correction, PFC) katı ve hücre dengeleme dönüştürücüleri ayrı yarı iletken gruplarıyla gerçekleştirilebilir. Her ilave kat güç yoğunluğu, sürücü, soğutma, koruma ve maliyet üzerinde yeni yük oluşturur. Tez kapsamında geliştirilen GT-tipi mimari, T-tipi dönüştürücünün anahtarlarını çalışma kipine göre yeniden görevlendirerek evirici, Vienna PFC ve aktif batarya yönetimi işlevlerini ortak bir donanımsal yapı üzerinde birleştirmeyi amaçlamaktadır. Böylece düşük güçteki anahtarlama ve harmonik problemi ile şarj--deşarj mimarisinin yeniden yapılandırılması, birbirinden kopuk iki çalışma olarak değil, T-tipi dönüştürücünün yapısal olanaklarını farklı ölçeklerde kullanan tek bir araştırma ekseni olarak ele alınmaktadır.

Belirlenen çerçevede tezin ana savı şudur: \emph{T-tipi dönüştürücünün yapısal ve anahtarlama özellikleri, enerji depolama sistemlerinde şarj--deşarj işlevlerini yeniden yapılandırmak ve düşük güçte verim--harmonik ödünleşimini iyileştirmek için birlikte kullanılabilir; Sin\textsuperscript{3} terimi ise bu ikinci hedefi destekleyen bir ileri besleme kompanzasyonudur.} Literatürde T-tipi eviricinin yüksek güç yoğunluğu, şebeke kontrolü, hafif yük dönüştürücü kipleri ve ileri besleme yolları ayrı ayrı incelenmiştir \citep{Blaabjerg2006,Wang2020,Chen2010,Lu2016}. Tezin odaklandığı araştırma boşluğu; T-tipi faz kolundaki düşük güç ters akım mekanizmasını, bunun verim--THD sonucu ile Sin\textsuperscript{3} tabanlı kompanzasyonunu ve aynı donanımın EDS işlevleri için yeniden yapılandırılmasını ortak bir mimari ve doğrulama zinciri içinde ele almaktır.

Tezin temel amacı, enerji depolama sistemleri için T-tipi dönüştürücü tabanlı yeniden yapılandırılabilir bir şarj--deşarj mimarisi geliştirmek ve aynı topolojinin şebekeye bağlı düşük güç çalışmasında görülen verim ile harmonik bozulma sorunlarını birlikte iyileştirecek yöntemleri ortaya koymaktır. Genel amaç doğrultusunda;

\begin{enumerate}
  \item T-tipi eviricinin düşük güçteki akım yollarını ve ters akım dolaşım mekanizmasını açıklamak,
  \item ters akımı azaltmaya yönelik nötr anahtar susturma yöntemini geliştirmek ve klasik anahtarlamayla karşılaştırmak,
  \item anahtarlama değişikliğinin üçüncü harmonik üzerindeki etkisini belirlemek ve Sin\textsuperscript{3} tabanlı ileri besleme kompanzasyonunu tanımlamak,
  \item kompanzasyon yolunu mevcut kapalı çevrim akım kontrolünden ayırarak sistem dinamiği ve kararlılık üzerindeki etkisini incelemek,
  \item T-tipi yarı iletkenleri farklı şarj, deşarj ve dengeleme kiplerinde yeniden kullanan GT-tipi mimariyi tanımlamak ve uygulanabilirliğini mevcut kanıtlarla değerlendirmek
\end{enumerate}
amaçlanmıştır.

Tez aşağıdaki araştırma sorularına yanıt aramaktadır:

\begin{enumerate}[label=AS\arabic*.]
  \item T-tipi eviricinin düşük güç çalışma bölgesinde ters akım dolaşımı hangi anahtarlama durumları ve akım yolları üzerinden oluşmakta; söz konusu dolaşım kayıpları ve akım harmoniklerini nasıl etkilemektedir?
  \item Nötr anahtarın koşula bağlı olarak susturulması, klasik anahtarlamaya göre düşük güç verimini hangi çalışma aralığında ve hangi harmonik maliyetle değiştirmektedir?
  \item Sin\textsuperscript{3} tabanlı ileri besleme kompanzasyonu, önerilen anahtarlamanın oluşturduğu üçüncü harmonik bileşenini ne ölçüde azaltmakta ve mevcut kapalı çevrim akım kontrolünün kararlılık marjlarıyla nasıl etkileşmektedir?
  \item T-tipi dönüştürücünün yarı iletkenleri, EDS'nin şarj--deşarj ve hücre dengeleme kiplerinde güvenli biçimde yeniden görevlendirilerek ayrı güç dönüşüm katlarının işlevleri ortak bir yapıda birleştirilebilir mi?
\end{enumerate}

Tezin özgün katkıları, birbirini tamamlayan üç düzeyde tanımlanmıştır:

\begin{enumerate}
  \item \textbf{Düşük güç anahtarlama yöntemi:} T-tipi eviricinin sıfır gerilim durumunda ortaya çıkan ters akım dolaşımını hedefleyen nötr anahtar susturma yöntemi geliştirilmiştir. Katkı; çalışma kiplerinin analizi, klasik ve önerilen akım yollarının karşılaştırılması, benzetim sonuçları ve erişilebilir deneysel kayıtlarla değerlendirilmektedir.
  \item \textbf{Sin\textsuperscript{3} tabanlı ileri besleme kompanzasyonu:} Yeni anahtarlama yönteminin üçüncü harmonik üzerindeki etkisini azaltmak üzere şebeke işaretinden anlık olarak üretilen bir kompanzasyon terimi önerilmiştir. Katkı, bağımsız bir kontrolcü olarak değil mevcut kapalı çevrime eklenen ileri besleme bileşeni olarak modellenmekte; harmonik performansın yanı sıra kapalı çevrim varsayımları ve kararlılık sınırlarıyla birlikte incelenmektedir.
  \item \textbf{GT-tipi yeniden yapılandırılabilir mimari:} Evirici, Vienna PFC ve aktif batarya yönetimi işlevlerinde aynı yarı iletkenlerin farklı görevlerle kullanılmasını sağlayan T-tipi tabanlı bir EDS mimarisi geliştirilmiştir. Katkı; çalışma kipleri, anahtarlama durumları, enerji akış yolları, güvenlik kısıtları ve mevcut benzetim/deney kanıtı üzerinden klasik ayrık mimariyle karşılaştırılmaktadır.
\end{enumerate}

Söz konusu katkılara ilişkin patent tescilleri teknolojik yenilik ve uygulanabilirlik bakımından tamamlayıcı çıktı niteliğindedir; bilimsel doğrulamanın yerine kullanılmamaktadır. Her katkının geçerliliği, tez içinde açıkça belirtilen model, benzetim, ölçüm ve karşılaştırma sınırları içerisinde değerlendirilmektedir.

Çalışma, şebekeye bağlı üç seviyeli T-tipi dönüştürücünün düşük güç davranışı ile bu topolojiden türetilen GT-tipi enerji depolama mimarisiyle sınırlandırılmıştır. Düşük güç incelemesinde anahtarlama durumları, ters akım, dönüştürücü verimi, şebeke akımı THD'si ve üçüncü harmonik temel performans ölçütleridir. Sin\textsuperscript{3} terimi, şebeke geriliminden üretilen ileri besleme kompanzasyonu olarak ele alınmakta; bağımsız bir kapalı çevrim kontrol yasası veya genel amaçlı aktif harmonik filtre yöntemi olduğu iddia edilmemektedir.

Doğrulama, tez kapsamında erişilebilen benzetim modelleri, gerçek zamanlı uygulama sonuçları ve deneysel kayıtlarla sınırlıdır. Sonuçlar; kullanılan DA bara gerilimi, şebeke koşulları, LCL filtre, anahtarlama frekansı, denetleyici ayarları, ölçüm düzeni ve incelenen güç seviyeleri için geçerlidir. Yeni deneyle doğrulanmamış çalışma bölgelerine ilişkin ifadeler genellenmemekte; benzetim ve deney sonuçları ayrı biçimde raporlanmaktadır. GT-tipi mimarinin maliyet, hacim ve eleman sayısı bakımından sağladığı yararlar, ayrıntılı yaşam döngüsü maliyeti veya seri üretim güvenilirlik çalışması yapılmadığı durumlarda ``potansiyel'' olarak ifade edilmektedir.

\eksikistek{Problem--yöntem--katkı akış şekli}{EDS $\rightarrow$ T-tipi dönüştürücü $\rightarrow$ düşük güçte ters akım $\rightarrow$ nötr anahtar susturma $\rightarrow$ Sin\textsuperscript{3} kompanzasyonu $\rightarrow$ GT-tipi mimari zincirini tek bakışta gösteren kavramsal bir şekil istiyoruz (teknik blok diyagramından farklı; TikZ veya MATLAB ile vektörel). Bölüm~4--7'nin nihai kanıtları kesinleştiğinde üretilip bu noktaya yerleştirilecektir.}

Tez sekiz bölümden oluşmaktadır. Birinci bölümde araştırma problemi genelden özele kurulmuş; amaç, araştırma soruları, katkılar ve sınırlar tanımlanmıştır. İkinci bölümde EDS güç dönüşüm mimarileri, T-tipi eviriciler, düşük güç verimi ve harmonik bozulma, akım kontrolü ve ileri besleme yaklaşımları incelenerek araştırma boşluğu belirlenmektedir. Üçüncü bölüm ortak sistem mimarisini, T-tipi çalışma kiplerini, LCL filtre ve şebeke modelini, kontrol yapısını, elektriksel parametreleri ve doğrulama yöntemini sunmaktadır.

Dördüncü bölümde düşük güç ters akım mekanizması çözümlenmekte ve nötr anahtar susturma yöntemi klasik anahtarlamayla karşılaştırılmaktadır. Beşinci bölüm Sin\textsuperscript{3} tabanlı ileri besleme kompanzasyonunu matematiksel olarak tanımlamakta; kapalı çevrim modelini, duyarlılık fonksiyonlarını ve kararlılık incelemesini vermektedir. Altıncı bölüm, benzetim ve erişilebilir deneysel sonuçları ortak bir test matrisi üzerinden değerlendirerek yöntemlerin düşük, orta ve mümkün olan anma gücü koşullarındaki performansını karşılaştırmaktadır. Yedinci bölüm GT-tipi yeniden yapılandırılabilir EDS mimarisini, çalışma kiplerini ve klasik mimariyle farklarını sunmaktadır. Sekizinci ve son bölümde araştırma soruları kanıt sınırları içinde yanıtlanmakta, özgün katkılar topluca değerlendirilmekte ve gelecek çalışmalar önerilmektedir.
